% GENERATED FILE. Edit canonical lessons, then run npm run generate:books.
% canonical-source-hash: fnv1a64:bf02f6a5d3b58343
% canonical-lessons: KA-C169-talme, KA-C169-asuye, KA-C169-ontitana, KA-C169-kutuhala, KA-C169-bhavane

\chapter{Patience, Jealousy, Loneliness, Curiosity, Feeling}
\label{ch:ka-patience-jealousy-loneliness-curiosity-feeling}
\hlchaptermodality{169}

\begin{chapteropening}
\textbf{By the end of this chapter:} \emph{I can say patience, jealousy, loneliness, curiosity and a feeling.}

Complete the last lesson of chapter 169: I can say patience, jealousy, loneliness, curiosity and a feeling.
\end{chapteropening}

\section[\texorpdfstring{\kn{ತಾಳ್ಮೆ} (tāḷme)}{ತಾಳ್ಮೆ (tāḷme)}]{\kn{ತಾಳ್ಮೆ} (tāḷme) — patience}

\label{lesson:KA-C169-talme}



\begin{quote}
Before the new one: say the Kannada for the thigh, then the Kannada for the palm.
\end{quote}

\subsection*{You'll want to know: \kn{ತಾಳ್ಮೆ}}
\textbf{\kn{ತಾಳ್ಮೆ}} — \emph{tāḷme} — \textquotedblleft{}patience\textquotedblright{}.

\begin{grammarlens}[title={What you've built}]
One more word for this chapter.
\end{grammarlens}

\subsection*{Your turn}
\begin{itemize}
\raggedright
  \item \emph{Say it:} \emph{tāḷme}
  \item \emph{Say it:} \emph{tāḷme}, once more
  \item \emph{Say it:} say \emph{tāḷme}
  \item \emph{Recall:} say the Kannada for the thigh, then the Kannada for the palm, then say \emph{tāḷme} again
\end{itemize}

\begin{tcolorbox}[breakable,colback=teal!4,colframe=teal!35!black,title={Before you move on}]
What does \kn{ತಾಳ್ಮೆ} mean? (\textquotedblleft{}Patience\textquotedblright{}.) Say it once more.
\end{tcolorbox}
\section[\texorpdfstring{\kn{ಅಸೂಯೆ} (asūye)}{ಅಸೂಯೆ (asūye)}]{\kn{ಅಸೂಯೆ} (asūye) — jealousy, envy}

\label{lesson:KA-C169-asuye}



\begin{quote}
Before the new one: say the Kannada for the palm, then the Kannada for patience.
\end{quote}

\subsection*{You'll want to know: \kn{ಅಸೂಯೆ}}
\textbf{\kn{ಅಸೂಯೆ}} — \emph{asūye} — \textquotedblleft{}jealousy, envy\textquotedblright{}.

\begin{grammarlens}[title={What you've built}]
One more word for this chapter.
\end{grammarlens}

\subsection*{Your turn}
\begin{itemize}
\raggedright
  \item \emph{Say it:} \emph{asūye}
  \item \emph{Say it:} \emph{asūye}, once more
  \item \emph{Say it:} say \emph{asūye}
  \item \emph{Recall:} say the Kannada for the palm, then the Kannada for patience, then say \emph{asūye} again
\end{itemize}

\begin{tcolorbox}[breakable,colback=teal!4,colframe=teal!35!black,title={Before you move on}]
What does \kn{ಅಸೂಯೆ} mean? (\textquotedblleft{}Jealousy, envy\textquotedblright{}.) Say it once more.
\end{tcolorbox}
\section[\texorpdfstring{\kn{ಒಂಟಿತನ} (oṇṭitana)}{ಒಂಟಿತನ (oṇṭitana)}]{\kn{ಒಂಟಿತನ} (oṇṭitana) — loneliness}

\label{lesson:KA-C169-ontitana}



\begin{quote}
Before the new one: say the Kannada for patience, then the Kannada for jealousy.
\end{quote}

\subsection*{You'll want to know: \kn{ಒಂಟಿತನ}}
\textbf{\kn{ಒಂಟಿತನ}} — \emph{oṇṭitana} — \textquotedblleft{}loneliness\textquotedblright{}.

\begin{grammarlens}[title={What you've built}]
One more word for this chapter.
\end{grammarlens}

\subsection*{Your turn}
\begin{itemize}
\raggedright
  \item \emph{Say it:} \emph{oṇṭitana}
  \item \emph{Say it:} \emph{oṇṭitana}, once more
  \item \emph{Say it:} say \emph{oṇṭitana}
  \item \emph{Recall:} say the Kannada for patience, then the Kannada for jealousy, then say \emph{oṇṭitana} again
\end{itemize}

\begin{tcolorbox}[breakable,colback=teal!4,colframe=teal!35!black,title={Before you move on}]
What does \kn{ಒಂಟಿತನ} mean? (\textquotedblleft{}Loneliness\textquotedblright{}.) Say it once more.
\end{tcolorbox}
\section[\texorpdfstring{\kn{ಕುತೂಹಲ} (kutūhala)}{ಕುತೂಹಲ (kutūhala)}]{\kn{ಕುತೂಹಲ} (kutūhala) — curiosity}

\label{lesson:KA-C169-kutuhala}



\begin{quote}
Before the new one: say the Kannada for jealousy, then the Kannada for loneliness.
\end{quote}

\subsection*{You'll want to know: \kn{ಕುತೂಹಲ}}
\textbf{\kn{ಕುತೂಹಲ}} — \emph{kutūhala} — \textquotedblleft{}curiosity\textquotedblright{}.

\begin{grammarlens}[title={What you've built}]
One more word for this chapter.
\end{grammarlens}

\subsection*{Your turn}
\begin{itemize}
\raggedright
  \item \emph{Say it:} \emph{kutūhala}
  \item \emph{Say it:} \emph{kutūhala}, once more
  \item \emph{Say it:} say \emph{kutūhala}
  \item \emph{Recall:} say the Kannada for jealousy, then the Kannada for loneliness, then say \emph{kutūhala} again
\end{itemize}

\begin{tcolorbox}[breakable,colback=teal!4,colframe=teal!35!black,title={Before you move on}]
What does \kn{ಕುತೂಹಲ} mean? (\textquotedblleft{}Curiosity\textquotedblright{}.) Say it once more.
\end{tcolorbox}
\section[\texorpdfstring{\kn{ಭಾವನೆ} (bhāvane)}{ಭಾವನೆ (bhāvane)}]{\kn{ಭಾವನೆ} (bhāvane) — a feeling, an emotion}

\label{lesson:KA-C169-bhavane}



\begin{quote}
Before the new one: say the Kannada for loneliness, then the Kannada for curiosity.
\end{quote}

\subsection*{You'll want to know: \kn{ಭಾವನೆ}}
\textbf{\kn{ಭಾವನೆ}} — \emph{bhāvane} — \textquotedblleft{}a feeling, an emotion\textquotedblright{}.

\begin{grammarlens}[title={What you've built}]
One more word for this chapter.
\end{grammarlens}

\subsection*{Your turn}
\begin{itemize}
\raggedright
  \item \emph{Say it:} \emph{bhāvane}
  \item \emph{Say it:} \emph{bhāvane}, once more
  \item \emph{Say it:} say \emph{bhāvane}
  \item \emph{Recall:} say the Kannada for loneliness, then the Kannada for curiosity, then say \emph{bhāvane} again
\end{itemize}

\begin{tcolorbox}[breakable,colback=teal!4,colframe=teal!35!black,title={Before you move on}]
What does \kn{ಭಾವನೆ} mean? (\textquotedblleft{}A feeling, an emotion\textquotedblright{}.) Say it once more.
\end{tcolorbox}
